\section{Conclusion}
\label{sec:conclusion}

We introduced a support-calibrated artifact-subspace diffusion model that
personalizes an estimated ocular propagation geometry rather than a clean-brain
prior or participant identity.  The support basis defines the rank-two latent
state and the observation-preserving reconstruction, making it impossible for
the implementation to ignore the operator while still producing the declared
output.

On 16 paired Klados source records, the diffusion point estimate improved clean
RRMSE utility over an information-matched deterministic estimator by
$0.0592$ (95\% paired-bootstrap interval $[0.0336,0.0879]$).  Its sample
variance also ranked errors more effectively than a three-seed deterministic
ensemble.  On 58 of 59 SGEYESUB stems, the frozen natural-EEG preservation
margins passed.  In contrast, matching calibration did not improve the
independent EOG endpoint over the population operator
($-0.00515$, $[-0.0113,0.00071]$) or the same-cell wrong operator
($0.00286$, $[-0.00303,0.00867]$).  The supported conclusion is therefore
limited: iterative diffusion is useful for the tested artifact-subspace target,
but a subject-specific increment from the present support estimator is not
established.

Klados records are not reliable participant identities and SGEYESUB was used in
development.  Moreover, the public EEGEyeNet release could not be retrieved
because its official Google Drive denied repeated participant-file access.
Independent participant-level confirmation remains necessary before presenting
the method as a validated subject-aware denoiser.
